\chapter{Three years, 1959--1962}
\label{ch:revolution}

\section{The question of the turn}

The single most argued question in modern Cuban history is whether Fidel Castro
was a communist in January 1959 who concealed it, or a radical nationalist whom
American reaction drove into the Soviet camp. Sixty years of scholarship has
not settled it and this book will not either, but the shape of the evidence is
worth stating, because Part III's treatment of the United States relationship
depends on which lesson one draws.

Against the concealment thesis: the July 26th Movement's published programme
was nationalist and constitutionalist; the Communist party opposed the
insurrection almost to the end and was distrusted by the guerrillas for it; the
first cabinet was liberal, headed by the judge Manuel Urrutia as president and
the lawyer Jos\'e Mir\'o Cardona as prime minister, with figures like Felipe
Pazos --- an orthodox economist --- at the central bank; and Castro spent April
1959 in the United States assuring audiences he was not a communist, at a
moment when saying so had no obvious tactical value he could not have obtained
more cheaply.

For it: Ra\'ul Castro and Che Guevara were Marxists well before 1959 and neither
was a marginal figure; the parallel structures that displaced the formal
cabinet were in place within months; the pace at which the liberal ministers
were removed --- Urrutia gone by July 1959, Pazos by November, Huber Matos
arrested in October for objecting to communist influence in the army and
sentenced to twenty years --- was too fast to be improvised; and Castro himself
said in 1961 that he had been a Marxist-Leninist all along, though he said many
things in 1961.

The most defensible reading is that neither pure version survives the record
\citep{farber2006}.
The revolution's leadership contained committed Marxists and committed
liberals, the leader was neither and was above all a nationalist with an
extraordinary instinct for where power lay, and the sequence of 1959--60 was a
genuine interaction in which each side's move made the other's next move
cheaper. The agrarian reform threatened American holdings; the American
response threatened the reform; the Soviet alternative made the response
survivable; taking it required the internal alliance that produced the party.
By the end of 1960 there was no path back for anybody, and the question of what
was in Castro's head in January 1959 had stopped mattering.

What matters for Part III is the second-order lesson: Cuba's sovereignty
problem was never solved, only re-addressed. The country moved from dependence
on one great power to dependence on another, and in 1990 discovered what that
had cost. A design for 2026 that does not think hard about the structure of
external dependence --- who buys, who lends, who supplies the fuel --- will
reproduce the same vulnerability with different flags.

\section{1959: the year of the two governments}

Formally, Cuba in 1959 had a cabinet, a president and a restored 1940
constitution as amended by a Fundamental Law. Actually it had the cabinet and
it had the Rebel Army, and where they disagreed the Rebel Army won. Castro took
the prime ministership in February. Urrutia was forced out in July after
criticising communist influence, on television, by Castro, in an act of
political theatre that established the method for the next decade.

Two things happened in 1959 that fixed everything after.

\subsection{The tribunals}

Batista's police and army had tortured and killed, largely with impunity, and
in January 1959 there was overwhelming popular demand for retribution and no
functioning court system with the legitimacy to deliver it. The revolutionary
government tried several hundred people before summary military tribunals and
executed them, most in the first four months. The number is disputed: figures
around 550 for 1959 are commonly cited and the Cuba Archive project's
compilations run considerably higher for the decade as a whole, while the
government has never published a count.\unv{} The trials were public, were
attended by crowds, and were in many individual cases against people who had
demonstrably committed the acts charged. They were also summary, appealed to no
independent authority, and in the case of the acquitted Batista airmen retried
to conviction on Castro's public insistence --- which destroyed, in the first
weeks, the principle that a verdict was final.

Both facts are true and they should be held together. A revolutionary
government facing genuine mass demand for justice against a state that had
tortured people is in a real dilemma, and doing nothing was not available. What
the tribunals established, however, was that in the new Cuba the law was an
instrument of the executive and not a check on it, and that precedent was never
subsequently reversed. The Huber Matos trial in October used the same machinery
against a revolutionary commander whose offence was political disagreement, ten
months later.

\subsection{The agrarian reform}

The First Agrarian Reform Law of 17 May 1959 capped individual holdings at 30
\emph{caballer\'ias}, about 402 hectares, with exceptions to about 1,340
hectares for high-productivity land; expropriated the excess with compensation
in twenty-year bonds at 4.5 per cent, valued at the owners' own declared tax
assessments; banned foreign landownership; and distributed land to tenants and
squatters in inalienable parcels of about 27 hectares. It created INRA, the
National Institute of Agrarian Reform, which within two years was a parallel
government running credit, distribution, housing and much of the economy.

The reform was, by Latin American standards of the period, not extreme --- the
cap was generous, compensation was offered, and comparable reforms were being
urged on the region by the United States itself two years later under the
Alliance for Progress. What made it a rupture was the valuation and the
currency. American owners had for decades declared their land at a fraction of
its market value to reduce tax; the law took them at their word, and offered
payment in Cuban bonds no one believed would be honoured. Washington demanded
prompt, adequate and effective compensation in cash. Havana declined. That
disagreement, more than any speech, is where the two countries separated.

\section{1960: the mechanics of the break}

The year runs as a sequence of moves in which each side's action made the
other's next one nearly automatic.

February: Anastas Mikoyan brings a Soviet trade delegation to Havana and signs
an agreement to buy Cuban sugar and supply crude oil on credit. May:
diplomatic relations with Moscow restored. June: the American- and
British-owned refineries in Cuba --- Esso, Texaco, Shell --- refuse, on
instruction, to refine Soviet crude; Cuba seizes them. July: the United States
cancels the balance of the Cuban sugar quota. Also July: the Soviet Union
announces it will buy the sugar. August: Cuba nationalises the American
utilities, telephone company, sugar mills and refineries. October: the United
States imposes an embargo on all exports to Cuba except food and medicine. Also
October: Cuba nationalises essentially the remainder of large enterprise ---
banks, some 380 large domestic firms, and, under the Urban Reform Law, all
rental housing, which was converted into ownership for the sitting tenants at
their existing rent paid over five to twenty years.

That last measure deserves more attention than it usually gets, because it is
still shaping Cuba. The Urban Reform Law made a majority of Cuban households
owners of their dwellings overnight and abolished the residential rental market
for six decades. It was massively popular; it is the single most durable
redistribution the revolution carried out; and it produced a housing stock in
which nobody had any incentive or legal mechanism to invest in maintenance, and
which has been falling down ever since. Havana loses buildings to collapse every
year. Chapter~\ref{ch:capital} has to deal with the consequence, which is that
Cuba has near-universal home ownership and a catastrophically decayed housing
stock, and that the two facts are the same fact.

January 1961: Washington breaks diplomatic relations. February 1962: the
embargo is extended to virtually all trade, under the Trading with the Enemy
Act, where it remains.

\section{The exodus begins}

Between 1959 and 1962 roughly 250,000 Cubans left, out of 6.6 million
\citep{eckstein2003}.\unv{}
They were disproportionately the professional and business class: the
expropriated, and then their doctors, lawyers, engineers and managers. Around
half of Cuba's physicians left in the first years --- roughly three thousand of
six thousand --- which is the origin of the revolution's medical policy, since
the state's response was to build medical schools at a rate no comparable
country attempted \citep{feinsilver1993}.\unv{}

Two features of this first wave matter later. It was legal and largely orderly,
by air to Miami, and it was understood by those leaving as temporary --- most
expected to return within months, which is why so much property was left in the
care of relatives, a fact that will make restitution claims in
Chapter~\ref{ch:capital} unusually complicated. And it included Operation Peter
Pan, in which some 14,000 unaccompanied children were flown to the United
States between 1960 and 1962 by a Catholic-organised programme, on the strength
of a rumour that the state intended to take custody of children. The rumour was
false. The programme was real, many of the families were separated for years,
and the episode poisoned the emotional register of the Cuban--American question
in a way that has never fully cleared.

\section{Playa Gir\'on and October 1962}

On 17 April 1961 about 1,400 Cuban exiles, trained, armed and transported by
the CIA, landed at the Bay of Pigs. The operation had been planned under
Eisenhower and inherited by Kennedy, who cancelled the second air strike; the
air cover failed, the expected internal rising did not occur --- the security
services had rounded up tens of thousands of suspected sympathisers in the
preceding days --- and the brigade was defeated in seventy-two hours. Around
1,200 were captured and later ransomed for food and medicine.

The invasion was the most consequential single American decision of the period,
and its consequence was the opposite of its intent. It converted the Cuban
government's security paranoia into a demonstrated fact, it justified permanent
mobilisation, it discredited the internal opposition by association with a
foreign landing, and it gave Castro the occasion --- on 16 April, at the funeral
of the victims of the preliminary air raid --- to declare the socialist
character of the revolution for the first time. It also made Cuba's request for
Soviet military guarantees reasonable rather than paranoid.

Eighteen months later, in October 1962, American reconnaissance found Soviet
medium-range ballistic missiles under construction in Cuba. The thirteen days
that followed were the closest the world has come to nuclear war. The
resolution --- Soviet withdrawal in exchange for an American undertaking not to
invade, and a quiet withdrawal of American missiles from Turkey --- was
negotiated between Moscow and Washington over Havana's head and against
Castro's furious objection; he had wanted the missiles to stay and had, at the
peak of the crisis, urged Khrushchev to consider a first strike if the United
States invaded.

Two legacies. The non-invasion understanding is, in practice, what has kept
Cuba's government alive for sixty years, and it was obtained by the Soviet
Union rather than by Cuba. And the experience of being traded over taught the
Cuban leadership a lesson about great-power patronage that it then declined to
act on for another twenty-eight years.

\section{Consolidation at home}

The same three years built the domestic machinery.

The literacy campaign of 1961 sent about 100,000 volunteers, most of them
teenagers, into the countryside for eight months. It reported teaching some
707,000 adults to read and reducing illiteracy from around 23 per cent to under
4.\unv{} The standard applied was elementary and the figure has been challenged
at the margins, but the campaign's reality is not seriously in doubt --- it is
visible in later cohort data --- and it remains the most impressive thing the
revolution did in its first decade. It was also, and simultaneously, an
instrument of political integration: the brigadistas carried a manual whose
reading exercises were about the revolution, and the campaign put a hundred
thousand urban adolescents under state discipline in the countryside for a
year.

In June 1961 private schools were nationalised, ending Catholic education. The
same year saw the founding of the Committees for the Defence of the Revolution,
a block-by-block network that combined vaccination drives, blood donation and
civil defence with neighbourhood surveillance --- one organisation doing both,
which is precisely why it worked and precisely why it was feared.
Independent press ceased to exist over 1960--61 through a mixture of
expropriation, printers' union takeovers and emigration. In March 1962 rationing
was introduced, the \emph{libreta}, intended as a temporary emergency
measure; it is still in force in 2026.

By 1962 the revolutionary organisations had been merged into the ORI, and its
first crisis --- the removal of the old communist An\'ibal Escalante in March
1962 for staffing the apparatus with his own people, denounced by Castro as
``sectarianism'' --- established who would be running the party. It was not the
communists.

\begin{ledger}{1959--1962}
\emph{Achieved.} Land redistributed to over 100,000 tenant families; rents
abolished and urban dwellings transferred to occupants; illiteracy effectively
eliminated in a single year; racial segregation of public and semi-public space
ended by decree and enforced; the rural half of the country brought inside the
state's reach for the first time; foreign ownership of the commanding heights
ended; and an invasion defeated.

\emph{Cost.} Summary executions in the hundreds, with no appeal and no
published count; the courts subordinated to the executive within weeks and
never restored; the independent press extinguished within two years; every
political organisation other than the governing one dissolved; a quarter of a
million people, including half the physicians, gone; private education
abolished; and the beginning of a security apparatus whose block-level reach
was without precedent in the hemisphere. Rationing introduced as a temporary
measure sixty-four years ago.
\end{ledger}
